\documentclass[english,screen,twoside]{noahnotes}

\begin{filecontents*}[overwrite]{references.bib}
@book{ogata2010,
  author    = {Katsuhiko Ogata},
  title     = {Modern Control Engineering},
  edition   = {5},
  publisher = {Prentice Hall},
  date      = {2010}
}
\end{filecontents*}

\addbibresource{references.bib}

\title{Control Systems Notes}
\author{Example Student}
\studentid{202600001}
\date{8 October 2026}

\documentsetup{
  course={Advanced Control Systems},
  course-code={MECH-500},
  course-short={Control},
  instructor={Ada Lovelace},
  semester={Autumn 2026},
  logo={../../figures/AU.pdf},
  toc-depth=2
}

\begin{document}

\frontmatter
\maketitle
\tableofcontents

\mainmatter

\lecture[State models]{1}{8 October 2026}{State-space models}
\label{lec:state-models}

\section{Definitions and notation}

Let \(\Vec{x}\) denote the state vector and \(\Mat{A}\) the state matrix.
A standard introduction is given by \textcite{ogata2010}.

\begin{definition}[State]\label{def:state}
The state is the smallest collection of variables that determines the future
behaviour of a system when its future input is known.
\end{definition}

\begin{example}[Second-order system]\label{ex:second-order}
A linear state-space model may be written as
\[
  \dot{\Vec{x}} = \Mat{A}\Vec{x} + \Mat{B}\Vec{u}.
\]
\end{example}

\begin{theorem}[State-transition matrix]\label{thm:stm}
For a continuous matrix \(\Mat{A}(t)\), a state-transition matrix exists.

\proofpart
This follows from existence and uniqueness for linear ordinary differential
equations.
\end{theorem}

As stated in \Mref{def:state}, the selected variables must determine the
future response. See also \cref{thm:stm}.

\backmatter
\printbibliography[heading=bibintoc,title={References}]

\end{document}
